% WRITE THE ABSTRACT IN THIS FILE.
% The abstract is a brief-by-design summary; per the analyst, it is a
% qualified exception to carrying its own PROVENANCE block, since the quoted
% sensitivity and luminosities are a plain cross-reference to
% Sec.~\ref{sec:results}, which carries the full provenance for them.
This note documents a measurement of
$R(D^{\pm}) = \mathcal{B}(B^{0} \to D^{-} \tau^{+} \nu_{\tau}) /
              \mathcal{B}(B^{0} \to D^{-} \ell^{+} \nu_{\ell})$
at Belle~II, using an inclusive-tagging approach in which the signal-side
$B$ is reconstructed as $B^{0} \to D^{-} \ell^{+}$ with
$D^{-} \to K^{+} \pi^{-} \pi^{-}$, and the tag-side $B$ is inferred from
rest-of-event information rather than from an exclusive tag reconstruction.
The signal yield in this inclusive approach is much larger than in the FEI method. However, there are difficult systematic and technical problems with the large $B\bar{B}$ backgrounds.
Signal and normalization yields are extracted from a two-dimensional binned
likelihood fit to the missing-mass squared and the scalar sum of the
$D$ and lepton center-of-mass momenta, performed simultaneously in a
$D$-mass signal region and a $D$-mass sideband. A preliminary Asimov-sensitivity study (Sec.~\ref{sec:results}), combining
the electron and muon channels, gives an expected relative uncertainty
of approximately $12\%$ on the $B^{0}\to D^{-}\tau^{+}\nu_{\tau}$
normalization for the current Run~1 ($365~\mathrm{fb}^{-1}$) and a partial Run~2 dataset ($125~\mathrm{fb}^{-1}$); this
sensitivity includes data statistics, MC statistics, and two of the normalization
systematics (continuum, and the $D^{*}\tau\nu$/$D^{**}\tau\nu$
feed-down); this is expected to change as further systematics are added to
the fit.

As of October 1, 2026, the signal region is blinded.
